\CatchFileBetweenTags{\AlphaInvVal}{constants.tex}{AlphaInvVal}

\begin{abstract}
Persistence Mechanics establishes that any system maintaining structural coherence against entropic decay must satisfy strict constraints of Finiteness, Conservation, and Causality, governed by six irreducible pillars. While intuitive at macroscopic scales, a fundamental question remains: do these architectural bounds apply to the substrate of reality itself? We test this by applying Persistence Mechanics to the zero-flux limit: the quantum vacuum. We find that satisfying these existential constraints strictly mandates the vacuum substrate to be positive-definite, even, and self-dual, uniquely isolating the manifold to a four-dimensional projection of the $E_8$ lattice. Evaluating the total geometric impedance of the six pillars on this substrate yields a dimensionless structural index $Z_{\text{geo}} = \num{\AlphaInvVal}\ldots$, containing zero continuous empirical parameters. This value structurally corresponds to the inverse fine-structure constant $\alpha^{-1}(0)$, agreeing with the kinematic measurement of Morel et al. (2020) within $0.58\,\sigma$ and CODATA (2022) within $1.68\,\sigma$, while predicting a \textit{Quantization Gap} where high-loop QED extractions systematically diverge. These results suggest that standard physical law represents the ultimate limiting case of structural persistence.
\end{abstract}